% !TeX root = ../../main.tex
\section{Introduction}
\label{sec:introduction}

Event cameras are being used in applications that require fast responses to visual changes or operation under a wide range of lighting conditions~\cite{gallego2022event}. Examples include visual localization and mapping in robotics~\cite{vidal2018ultimate}, steering prediction for autonomous driving~\cite{maqueda2018steering}, and eye tracking for near-eye displays~\cite{angelopoulos2021gaze}. These applications motivate both the development of event sensors and methods for processing their output.

Unlike a conventional camera that captures complete images at a fixed frame rate, a dynamic vision sensor (DVS) reports brightness changes independently at each pixel. When the change in log intensity reaches a threshold, the pixel emits an event containing its coordinates, timestamp, and the polarity of the change~\cite{lichtsteiner2008dvs}. The output is therefore a stream of events, commonly represented as $(x,y,t,p)$, rather than a sequence of intensity images.

This sensing technology provides high temporal resolution, high dynamic range, and low-latency processing of brightness changes~\cite{gallego2022event}. It also avoids the exposure-time integration responsible for conventional motion blur, although sensor bandwidth and latency still limit the response to fast motion~\cite{hu2021v2e}. These properties come with a different data format: events arrive at irregular times and do not directly form the dense images expected by conventional vision algorithms. Learning methods address this by processing event streams directly or converting them into representations suitable for various neural networks.

Many of these methods use supervised learning and therefore require labeled training data. Collecting real events and obtaining accurate labels can be difficult, particularly for fast motion or conditions where a conventional camera cannot provide reliable ground truth. This limits the available training data and motivates the use of synthetic events~\cite{hu2021v2e,lin2022dvsvoltmeter}.

Event-camera simulators make it possible to generate such data from existing videos, rendered scenes, or their 3D representations. For example, v2e converts intensity frames into events using a pixel model and can interpolate additional frames to improve temporal sampling~\cite{hu2021v2e}. DVS-Voltmeter models event generation through a stochastic process~\cite{lin2022dvsvoltmeter}, PECS includes the optical path and scene rendering~\cite{han2024pecs}, and V2CE learns a video-to-event conversion with continuous timestamp inference~\cite{zhang2024v2ce}. These approaches address different parts of the simulation process, so their outputs can differ in both spatial and temporal structure.

However, reproducing the response of a real sensor remains difficult. Work on physically realistic pixel models continues to address the effects of sensor circuitry, illumination, and noise~\cite{graca2025sensormodel}. Simulator evaluations commonly combine visual inspection with performance on selected downstream tasks~\cite{hu2021v2e,lin2022dvsvoltmeter}. Direct comparisons have also been proposed: PECS uses spatiotemporal event distances~\cite{han2024pecs}, while Event Quality Score (EQS) compares features extracted by a pretrained network~\cite{chanda2025eqs}. These measurements provide different views of the real--simulated gap, and their interpretation depends on the representation and comparison procedure.

Downstream performance shows how well synthetic data serve a particular task, but does not directly measure how closely they reproduce the information collected from real sensors. A direct comparison with recordings of the same stimulus would provide a separate assessment of simulator output and how it changes with parameter settings. To support such an assessment, the recordings must be comparable, and the resulting distances need a reference that helps explain their magnitude and sensitivity.

We propose a measurement framework for comparing real and simulated event streams under controlled conditions and examining how the resulting distances should be interpreted. The framework combines calibrated, paired recordings with maximum mean discrepancy (MMD) and sliced Wasserstein distance (SWD), alongside a Chamfer-distance baseline. We use these established methods to measure differences in event structure and examine their responses through controlled tests. The contribution is the data collection and comparison procedure, together with the experimental evidence supporting its use under specified configurations.

To obtain comparable data, we record a rotating optical chopper with an event camera and a frame camera. The frame sequence supplies input to video-based simulators, while calibration establishes spatial and temporal alignment with the real events. The chopper provides a simple stimulus with adjustable rotation speed and repeated patterns under nominally identical conditions. Its known geometry also provides a reference for refining spatial alignment.

This periodic stimulus supports several complementary tests. Random subsets of the same recording characterize distances caused by finite sampling, while observations from different rotations provide a real--real comparison. Controlled offsets establish a ruler relating measured distance to displacement of the stimulus, and modifications to the recorded events isolate sensitivity to particular changes. These references then support interpretation of real--simulated comparisons. The measurements remain dependent on choices such as the space--time scale, kernel bandwidth, and comparison window; the tests examine what can be concluded under those choices.

The remainder of the paper first introduces the sensor model and the comparison methods. The methodology then describes data acquisition, calibration, event representation, and the evaluation protocol. The evaluation presents the controlled tests and simulator comparisons, followed by a discussion of their implications and limitations.
